\chapter{Eigene Fortsetzung: ein gemeinsames kontrolliertes Mikrolabor}
\label{ch:micro14}
\begin{bild}{Aktualisierung v1.5: zwei gültige Präparationsverträge}
Der unten konstruierte exakt selektive Spektralfilter bleibt gültig unter seinen stärkeren Kontrollannahmen. Kapitel~\ref{ch:execution15} ergänzt eine einfachere Alternative: wiederholte Stern-Records liefern kontrollierte Näherungspräparation ohne kontrolliertes $e^{-iH\tau}$ und ohne 13-Faktor-Filter. Endliche Rundenzahl, Ausbeute und Controllerzeiten werden bezahlt.
\end{bild}
\section{Vorzeichen aus E8 und die noch fehlende Quellbasis}
Der geordnete Standard-Gittercocycle $\varepsilon(m,n)=(-1)^{m^TBn}$ wurde unabhängig aus einer einfachen E8-Basis aufgebaut. $B+B^T=G\pmod2$ liefert den Kommutatorfaktor, Bilinearität die Cocycle-Identität. Alle 240 Wurzelkoordinaten sind ganzzahlig. Auf den 40 Clebsch-Kanten bestätigt die exakte F2-Elimination Rang 45 für die Antisymmetriegauge der Ortsfarben und Rang 285 für die zusätzliche kantenlokale Vermittlergauge. Für die geteilte Bank existiert in der untersuchten Klasse diagonaler Vorzeichen keine Lösung.

Dies bestätigt die endliche Vorzeichenrechnung von N9. Es ist kein Beweis gegen alle komplexen oder nichtdiagonalen Adapter. Vor allem ist ein korrekt gewählter Standard-Gittercocycle noch nicht sein phasentreuer Transport zur tatsächlichen TFPT-Clock-/Quellbasis. Genau diese Schnittstelle bleibt offen.

\section{Ein allordentlicher CAR-/Tensoradapter im Matchingsektor}
Für ausschließlich kantenlokale Vermittler ohne Hopping sind die lokalen Ladungen
\begin{equation}
Q_v=n_f(v)+\sum_{e\ni v}n_e
\end{equation}
erhalten. Aus der vollständig belegten Nullbosonquelle folgt $Q_v=1$. Besetzte Bosonen bilden deshalb ein Matching: Jede Kante trägt höchstens einen, und zwei besetzte Kanten teilen keinen Endpunkt.

Wähle eine feste Ortsordnung und die konstante Kantenphase $(-1)^{j-i-1}$ für $i<j$. Ein CAR-Paarabbau unterscheidet sich dann vom harten Tensorabbau durch $(-1)^h$, wobei $h$ die Zahl der schon leeren Orte zwischen $i$ und $j$ zählt. Jedes alte Paar trägt null oder zwei zu $h$ bei, außer wenn seine Kante die neue Kante kreuzt. Definiere daher
\begin{equation}
G(M)=(-1)^{\#\{\text{Kreuzungen besetzter Kanten im Matching }M\}}.
\end{equation}
Für jeden Übergang $M\to M'$ gilt genau das relative Vorzeichen $G(M')G(M)$. Die Adjungierten folgen durch Rücknahme. Somit sind die \emph{vollständigen} CAR- und Tensorhamiltonoperatoren in diesem erreichbaren Sektor durch eine diagonale unitäre Gauge verbunden, nicht nur bis vierter Ordnung.

Der vollständige K4-Matchingraum hat 940 Zustände: 256 ohne Vermittler, 576 mit einem und 108 mit zwei. Alle 1584 Erzeugungsübergänge wurden exakt geprüft, darunter 144 nichttriviale relative Vorzeichen. Das Entfernen der Kreuzungszeichen fällt als Negativmutant durch.

\begin{grenze}{Der Scope gehört zum Satz}
Der Satz setzt lokale Kantenmoden, die feste Belegung $Q_v=1$, kein Vermittlerhopping und die erklärte Kantenphasenkalibrierung voraus. Er identifiziert nicht die geteilte Bank, andere Ladungssektoren oder den nativen Clock-Adapter. Eine spätere Transporterweiterung benötigt einen neuen Vergleich.
\end{grenze}

\section{Der vollständige Stern, nicht nur sein niedriger Schatten}
Die Kopplung $M$ zwischen 256 Materie- und 288 Vermittlerzuständen wurde neu aufgebaut. Sie erfüllt $M^\dagger M=6I-2G_\star$. Der volle Hamiltonoperator ist
\begin{equation}
H_\star=\begin{pmatrix}0&tM^\dagger\\tM&\Delta I_{288}\end{pmatrix}.
\end{equation}
Seine 14 verschiedenen Energien sind sieben niedrige Werte
\begin{equation}
E_-(g)=\frac{\Delta-\sqrt{\Delta^2+4t^2(6-2g)}}2,
\qquad g=0,\tfrac12,1,\tfrac32,2,\tfrac52,3,
\end{equation}
sechs obere $E_+(g)=\Delta-E_-(g)$ für $g<3$ und $\Delta$ mit Multiplizität 67. Die niedrigen Multiplizitäten sind $1,30,45,40,15,90,35$; die sechs oberen wiederholen die ersten sechs. Alle 544 numerisch berechneten Eigenwerte stimmen mit dieser Zuordnung überein. Zielzustand ist der eindeutige niedrigste angekleidete Zustand $\Omega_d$.

\section{Ein exakter endlicher Filter aus wirklichen Entwicklungszeiten}
Für jeden anderen Energiewert $E_j$ setze
\begin{equation}
\tau_j=\frac{\pi\hbar}{E_j-E_0},\qquad
A_j(H)=\frac{I+e^{-i\tau_j(H-E_0)/\hbar}}2.
\end{equation}
$A_j$ ist der akzeptierte Krausoperator eines Hadamardtests mit kontrollierter Entwicklung. Alle Faktoren erhalten den Zielzustand, und der $j$-te Faktor löscht den Eigenraum zu $E_j$. Weil sie Funktionen desselben $H$ sind, gilt auf dem ganzen Raum
\begin{equation}\boxed{\prod_{j=1}^{13}A_j(H_\star)=P_{\Omega_d}.}
\end{equation}
Es ist kein gemeinsames kommensurables Wiederkehrintervall nötig. Die vollständige 544D-Matrixrechnung bestätigt den Projektor; ein Produkt nur über die sechs anderen niedrigen Energien fällt am oberen und dunklen Sektor durch. Damit ist die Lücke zwischen einem algebraischen Filterpolynom und einem deklarierten kontrollierten Zeitprotokoll konkret verkleinert.

Bei $\Delta=1$, $t=1/20$ beträgt die Summe der 13 Zeiten
\begin{equation}
\sum_j\tau_j=3172{,}829634\,\hbar/\Delta.
\end{equation}
Die größte Einzelzeit ist $1290{,}728425\,\hbar/\Delta$. Benötigt werden bekannte Spektralwerte dieses kleinen Modells, genaue Phasen und Zeiten sowie kontrolliertes $H$. Das ist ein kontinuierlicher Kontrollvertrag, keine behauptete endliche exakte Clifford+T-Synthese.

Bei gemessenen Filtern genügt ein immer wieder frisch zurückgesetztes Bit; die vollständige Erfolgsbedingung umfasst 13 Ergebnisse. Kohärentes Markieren braucht hier 13 Hilfsbits: berechnen, den All-null-Record phasieren, zurückrechnen. Weil der akzeptierte Zweig genau der Projektor ist, liefert dies eine selektive Phase auf $\Omega_d$ und setzt die Hilfsbits zurück.

Eine Zielphase benötigt 26 kontrollierte Entwicklungen. Zwei angepasste Verstärkungsschritte benötigen 52 plus die beiden Eingangsphasen. Für die bereits berechnete Überlappung $a=w/6$ bleibt die Phase $1{,}734110745303\ldots$ korrekt. Bei sonst idealen Gattern begrenzt Duhamel/Teleskopieren den Zustandsfehler durch $4\sum_j\tau_j(\delta H+\delta E_0)/\hbar$. Für Infidelität $10^{-6}$ genügt konservativ $\delta H+\delta E_0\le7{,}88\cdot10^{-8}\Delta$. Zeit- und weitere Gatterfehler kommen hinzu.

\section{Noch einfacher: die Hülle messen und die Ausbeute mitführen}
Eine ideale kohärente Entkleidung ist für erfolgreiche Zustandspräparation nicht nötig. Nach dem Filter wird die Vermittlerbelegung gemessen. Im leeren Zweig bleibt exakt nacktes $\Omega$. Mit
\begin{equation}
w=\frac{1+\Delta/\sqrt{\Delta^2+24t^2}}2=0{,}985642931179\ldots
\end{equation}
gilt auf nackter Materie
\begin{equation}\boxed{
P_{\mathrm{bare}}P_{\Omega_d}P_{\mathrm{bare}}=wP_\Omega.}
\end{equation}
Dies ist der gesamte akzeptierte Krausoperator und wurde als volle Matrix geprüft. Aus dem elementaren Zweisingulett $\chi$ mit $|\langle\Omega,\chi\rangle|^2=1/6$ beträgt die Ausbeute $w^2/6$. Fehlgeschlagene Recordzweige sind nicht verschwunden, sondern ausdrücklich zurückzusetzen oder als Fehlschläge zu zählen.

Der resonante Record lässt sich dabei vereinfachen. Für den realen resonanten Volltransfer
\begin{equation}
U_{\mathrm{res}}=\begin{pmatrix}P_+&-iW^\dagger\\-iW&0\end{pmatrix}
\end{equation}
gilt auf allen 44 lokalen Dimensionen
\begin{equation}
(U_{\mathrm{res}}^\dagger\otimes I)Q(U_{\mathrm{res}}\otimes I)
=\bigl(P_+\otimes I+P_-\otimes X\bigr)\oplus I_{12}.
\end{equation}
Die zusätzlichen Viertelphasen werden nur für die spezielle Involution $U_0$ gebraucht, nicht für diese Record-Makrooperation. Beide Recordausgänge und die Rückkehr des Vermittlers sind eingeschlossen.

\section{Ein vollständiger neuer Präparations- und Echovertrag}
Das neue Protokoll beginnt in nacktem $\chi$, benutzt den mikroskopischen Sternfilter samt Leerbelegungsherald, danach den lokalen Dreierzyklus $C$, zweimal den resonanten Record, den Rücktick und \emph{denselben} mikroskopischen Filter zur Schlussprüfung. Der Echoanteil wurde im vollständigen 256D-Materieraum gerechnet; wiederholte Farben werden nicht herausprojiziert.

\begin{center}\begin{tabular}{lrr}
\toprule&Record behalten&Record frisch\\\midrule
Präparation erfolgreich&0{,}161915331297&gleich\\
Schlusserfolg nach Präparation&0{,}971491987782&0{,}516105118509\\
Gesamterfolg je Anfangsversuch&0{,}157299447054&0{,}083565331248\\\bottomrule
\end{tabular}\end{center}
Die exakten Formeln sind $w^2/6$ für Präparation, $w^2$ beziehungsweise $17w^2/32$ für die Schlussprüfung und $w^4/6$ beziehungsweise $17w^4/192$ für den Gesamterfolg. Frisch/behalten bleibt $17/32$. Die abweichenden Rohwerte gegenüber dem älteren $\ket{0123}$-Protokoll sind kein Widerspruch: Eingang, Filter und bezahlte Ausbeute sind ausdrücklich verschieden.

Bei Wiederholung mit Reset benötigt Präparation im Mittel $6{,}1761$ Versuche. Ein erfolgreicher Start-und-Schluss-Pfad umfasst 26 schwache kontrollierte Entwicklungen mit Gesamtzeit $6345{,}659268\hbar/\Delta$, zwei zusätzliche Leerbelegungsheralds und die Recordpulse. Diese endliche gemeinsame Modellrechnung ist geschlossen; die native Verfügbarkeit des Steuervertrags ist es nicht.

\section{Feedback: Attraktor, Fehlerboden und unabhängige Zellen}
Der neue N11-Kanal
\begin{equation}
\mathcal E(\rho)=K\rho K^\dagger+
\Tr[(I-K^\dagger K)\rho]\rho_0
\end{equation}
ist CPTP, mit $K=P^-_{03}P^-_{02}P^-_{01}$ und $\rho_0=\ket{0123}\bra{0123}$. Wegen $K\Omega=K^\dagger\Omega=\Omega$ und $K^\dagger K\le P_\Omega+\beta(I-P_\Omega)$ folgt
\begin{equation}
1-F(\mathcal E^m\rho)\le r^m(1-F\rho),\qquad
r=1-\frac{1-\beta}{24}=0{,}97542071045\ldots.
\end{equation}
Damit ist $P_\Omega$ der einzige stationäre Zustand. 556 Zyklen garantieren $10^{-6}$ Infidelität. Der unabhängige 256D-Lauf kontrolliert die Gewichtsidentität und Kontraktion in jedem der 556 Schritte. Dies ist nicht der ungelesene unitale Recordkanal mit großer Fixalgebra.

Eigene Erweiterung: Bei höchstens $\varepsilon$ Fehler pro Zyklus in halber Diamantnorm gilt
\begin{equation}
1-F_m\le r^m(1-F_0)+\frac{\varepsilon(1-r^m)}{1-r}.
\end{equation}
Der konservative Fehlerboden beträgt $40{,}6847\varepsilon$. Für $N$ unabhängig kontrollierte Zellen liefert der Vereinigungsbound sogar bei anfänglich verschränkten Eingängen $1-F_{\mathrm{global}}\le Nr^m$. Bei $N=4096$ genügen 890 Zyklen für global $10^{-6}$; dies ist kein wechselwirkender Zustandssatz.

Reset exportiert Information: Bis zu acht Farbbits je zurückgesetzter Zelle und drei Austauschrecords je Zyklus sind in dieser Ausführung zu berücksichtigen. Energie- und Entropiekosten hängen von Umgebung und Löschprotokoll ab. Ein geschlossenes endliches unitäres System erhält keine kostenlose reine Attraktion.
